\section{System framework}
\label{sec:system}

We propose the system flow shown in in Figure \ref{fig:sys_flow}.
Either the human supervisor or an automated object recognizer labels the search space to initialize a shared mental model.
With the shared mental model, a parse converts a verbal command from the human into a sub-task abstraction for the robot.
In the path planning problem, this abstraction is encoded as a mathematical model of the optimization problem.
Each adverb in the task grammar is associated with a different planning objective.
The path planner generates the set of Pareto optimal solutions and the human selects one that matches his/her intent. 
After a solution is found, it is sent to a robot to execute.

\begin{figure}[tbph]
\centering
\includegraphics[width=0.5\linewidth]{./images/sys_flow}
\caption{The process of a semantic-based path planning.}
\label{fig:sys_flow}
\end{figure}

\begin{figure}
\subfigure[Planned path from hexagonal map.]{
\centering
\includegraphics[width=0.28\linewidth]{./images/min_dist2}
\label{fig:plan}
}
\subfigure[Task execution monitor.]{
\centering
\includegraphics[width=0.29\linewidth]{./images/planMonitor3}
\label{fig:stage}
} 
\subfigure[Task execution in Gazebo simulator.]{
\centering
\includegraphics[width=0.37\linewidth]{./images/gazebo2}
\label{fig:stage-sim}
}
\caption{Simulation with the Gazebo simulator.}
\label{fig:simulation}
\end{figure}

We are currently developing tools to test our approach.
The map of a workspace is firstly labeled by a supervisor.
A labeled map is shown in Figure \ref{fig:Label}.
We assign semantic IDs to different sub-regions by using ``indoor'' and ``outdoor'' labels.
We label the doors and windows of several regions, which are frequently used in the verbal commands of a cordon and search task.
Moreover, we label some significant objects on the ground to facilitate localizing the positions of task execution.
Within the labeled map, a verb command, ``go quickly to location BD-2'', is read and parsed into multi-objectives and constraints by the task grammar.
A path planner finds an ``optimal'' solution through interactive multi-objective optimization.
The planning is currently on a discretized map, which is a hexagonal description of the test environment, which is shown in Figure \ref{fig:plan}.
Thus, the format of a planned path is a sequence of hexagonal cells.
This path is then converted into a sequence of waypoints, which is shown in Figure \ref{fig:stage}.
We test the task execution in the Gazebo simulator by using a virtual turtlebot, which is illustrated in Figure \ref{fig:stage-sim}.
It is easier to construct a testing urban environment in the Gazebo simulator, and the implementation can be easily migrated to a real turtlebot.
Figure \ref{fig:stage} shows a task execution monitor, which provides the progress of task execution.
A execution monitor GUI, shown in Figure \ref{fig:stage}, displays the sequence of waypoints of a planned path and continuously receives the position updates from the virtual robot, which is used to check how the task is executed.
The blue points are reached waypoints and the green points are the waypoints to be reached.
The process is shown in Figure \ref{fig:simulation}.
